% !TEX root = ../main.tex

\section{Related Work}
\label{sec:related_work}

\textbf{CNN encoder-decoder saliency models.} Our starting point is the VGG-based encoder-decoder of Wang and Shen~\cite{wang2018deep}, which we replicate as our first prototype (\secref{subsec:beta-v10}); a closely related design is Ghariba~et~al.'s five-conv/five-deconv VGG-16 encoder-decoder~\cite{ghariba2019visual}. Both rely on a pretrained encoder and a single reconstruction-style loss. Our from-scratch replica, trained with plain MSE, produced over-smoothed maps (\secref{subsec:beta-v10}); this gap motivated our departure from a single pixel-wise loss toward a deep-supervised, multi-term saliency loss (\secref{subsec:loss11}).

\textbf{Transfer learning from object recognition.} K\"ummerer~et~al.~\cite{kummerer2016deepgaze} show that VGG-19 features pretrained for object recognition, used \textit{without fine-tuning}, already predict fixations competitively, direct evidence for the hypothesis in \secref{subsec:litreview}. Their result is demonstrated only for convolutional backbones; in v2.0 (\secref{sec:model20}) we test whether it extends to a pretrained Swin Transformer~\cite{liu2021swin}, substituted for our convolutional encoder while the rest of the pipeline is held fixed.

\textbf{Transformer-based saliency.} TranSalNet~\cite{lou2022transalnet} augments a CNN backbone with bespoke, jointly-trained transformer fusion blocks, achieving state-of-the-art results on MIT300/SALICON at the cost of extra architectural and training complexity. v2.0 tests a cheaper alternative: an off-the-shelf pretrained Transformer \textit{encoder} paired with our existing lightweight decoders, which better fits our hardware budget (\secref{subsec:env}) and isolates the effect of the encoder alone.

\textbf{Evaluation methodology.} Yan~et~al.~\cite{yan2022review} and Bylinskii~et~al.~\cite{bylinskii2019different} argue that no single saliency metric suffices, since location- and distribution-based metrics capture different properties. We follow this by reporting three complementary distribution-based metrics (\secref{subsec:metrics}) and by testing cross-dataset on MIT1003~\cite{judd2009learning} rather than only on held-out SALICON data.

In short, we adopt the encoder-decoder topology and transfer-learning premise of~\cite{wang2018deep,ghariba2019visual,kummerer2016deepgaze}, replace the reconstruction loss with a deep-supervised saliency-specific loss, and test the attention-based direction of~\cite{lou2022transalnet} as a minimal, controlled encoder substitution rather than a bespoke fusion architecture.
